\appendix
\section{Reproducibility}
\label{app:repro}

\paragraph{Source code.} The complete pipeline, the six executed Jupyter
notebooks, all 44 result tables, all 33 figures and the LaTeX source of this
report are publicly available in the project repository:
\href{https://github.com/eqflashheqr-cmyk/DataMining}{\texttt{https://github.com/eqflashheqr-cmyk/DataMining}}
(released under the MIT licence). The raw dataset is not committed because of
its size --- it is the public ULB ``Credit Card Fraud Detection'' dataset on
Kaggle --- and neither are the trained \texttt{.joblib} models; every artefact
in this report is regenerated by the scripts in \texttt{src/}, and the
notebook pipeline is given in Table~\ref{tab:notebooks} below.

All experiments were executed with Python 3.11.9 on Windows. The numerical
core of the dependency set is pinned exactly in \texttt{requirements.txt}
(pandas 2.3.3, numpy 2.2.6, scipy 1.17.1, scikit-learn 1.9.1,
imbalanced-learn 0.14.2, xgboost 3.2.0, shap 0.51.0, joblib 1.6.0,
matplotlib 3.10.9, seaborn 0.13.2, plotly 7.1.0); only the notebook tooling
(jupyter, notebook) is left as a lower bound, since it does not affect any
reported number. Every figure and table in this report can therefore be
regenerated from the pinned versions alone. The global random
seed is fixed at \textbf{42} in \texttt{src/config.py} and applied to the
stratified split, every estimator, \texttt{GridSearchCV} and
\texttt{permutation\_importance}. The stratified partition is computed once and
persisted to \texttt{artifacts/models/split\_with\_features.joblib}; all
subsequent notebooks load this file rather than re-splitting, which guarantees
that every reported result refers to an identical partition.

\begin{table}[h!]
\centering
\caption{Notebook pipeline. Each notebook is executed in order and reads only
artefacts produced by its predecessors.}
\label{tab:notebooks}
\scriptsize
\setlength{\tabcolsep}{3pt}
\begin{tabular}{p{4.7cm}p{4.5cm}p{3.6cm}}
\toprule
\textbf{Notebook} & \textbf{Purpose} & \textbf{Runtime} \\
\midrule
\texttt{00\_setup\_and\_split.ipynb} & Load, clean, stratified split, feature engineering & $<$ 1 min \\
\texttt{01\_preprocessing.ipynb} & Scaler and feature-set ablation & $\sim$ 10 min \\
\texttt{02\_imbalance\_benchmark.ipynb} & 27-configuration benchmark & $\sim$ 73 min \\
\texttt{03\_model\_training.ipynb} & Final model selection and training & $\sim$ 48 min \\
\texttt{04\_evaluation\_shap.ipynb} & Evaluation, threshold tuning, SHAP & $\sim$ 2 min \\
\texttt{05\_cost\_analysis.ipynb} & Cost model, business questions & $<$ 1 min \\
\bottomrule
\end{tabular}
\end{table}

The deployed pipeline is \texttt{final\_model.joblib} (in
\texttt{artifacts/models/}) with its metadata in
\texttt{final\_model\_metadata.joblib}; the directory \texttt{artifacts/tables/}
holds 44 result tables and \texttt{artifacts/charts/} 33 figures.
Three implementation subtleties were resolved during the work and are recorded
here because they affect reproducibility. First, because the split has already
engineered the five new columns, the \texttt{ColumnTransformer} is declared
with \texttt{remainder="drop"} and an explicit list of the columns to
transform, so that any column not named is discarded deliberately rather than
passed through by accident. Second, a pipeline
containing a sampler must be built from \texttt{imblearn.pipeline.Pipeline},
not \texttt{sklearn.pipeline.Pipeline}, otherwise the intermediate step fails the
transformer contract. Third, the label arrays returned by the split retain the
\emph{original} CSV row indices, which are not contiguous, so positional access
must use \texttt{.iloc[]} rather than \texttt{[]}.

\paragraph{Preprocessing detail.} The \texttt{ColumnTransformer} in the final
pipeline declares two explicit transformers: a \texttt{StandardScaler} applied
to \texttt{Time} and \texttt{Amount}, and a \texttt{pass-through} applied to the
remaining 33 columns (\texttt{V1}--\texttt{V28} and the five engineered
features). The scaler is therefore fitted inside the pipeline, on training data
only. This design was chosen after comparing \texttt{StandardScaler},
\texttt{RobustScaler} and no scaling: scaling affects the linear models
strongly and the tree ensembles negligibly, so we retain
\texttt{StandardScaler} as the simplest option that is correct for both.

\paragraph{Daily-rate normalisation.} All per-day quantities reported in
Section~\ref{sec:cost} are obtained by multiplying the corresponding test-set
count by $5/3$. The derivation is: the cleaned dataset contains 283{,}726
transactions spanning approximately 48 hours, i.e.\ 141{,}863 transactions per
day; the test set contains 85{,}118 rows, i.e.\ $85{,}118 / 141{,}863 = 0.6$
days of transaction stream. The factor $1/0.6 = 5/3$ therefore converts a
test-set count into a daily rate. This normalisation was applied uniformly to
workload, cost and staffing figures, and the corresponding tables are archived
in \texttt{artifacts/tables/s10\_*.csv} together with the script
\texttt{src/s10\_policy\_corrected.py} that regenerates them.
